\subsection{The actual rank-four triangular alternatives}
\label{subsec:sr-four-actual-triangular-alternatives}

This subsection concerns the original rational representation,
before restriction to the real tower. Its root coordinates are
\[
 v_0=(1,1),\quad v_1=(0,1),\quad v_2=(1,0),\quad v_3=(0,0),
 \qquad w_{ij}=v_i-v_j.
\]
The two sectors are \(A=\{0,2\}\) and \(B=\{1,3\}\).
All coefficient extensions below are finite on a complete DVR
model. The arguments therefore apply both to characteristic zero
and to characteristic \(\ell\).

\paragraph{Integral coordinate projectors.}
Let \(\mathcal C=\operatorname{span}_{\mathcal V}\{r(g)\}\),
where \(\mathcal V\) is the valuation ring of an actual
normalized point. This span is a subalgebra and is stable under
symplectic adjoint: \(r(g)^\#=\nu(g)r(g^{-1})\).
The actual normalization supplies elements \(\varphi,g_*\) with
\(X=r(\varphi)\) block diagonal on \(A,B\), disjoint residual
spectra on those blocks, and
\[
 H=r(g_*)+\chi_0(g_*)r(g_*^{-1}),\qquad H_{02}=H_{20}=0.
 \tag{TA1}
\]
Here \(\chi_0=\chi_2\chi_1\), \(\chi_3=1\), and the residual
values of \(\chi_1,\chi_2\) at \(g_*\) differ from \(1\).
These are precisely the relations imposed by the original
normalization, rather than an assumption that an arbitrary
regular element is diagonal.

Let \(p_B\) be the characteristic polynomial of the \(B\)-block
of \(X\), and set \(Y=p_B(X)\).
Its \(B\)-block is zero and the determinant \(\Delta\) of its
\(A\)-block is a unit: its reduction is the product of all four
cross-block eigenvalue differences.
Cayley--Hamilton in the \(A\)-block gives
\[
 E_A=\frac{\operatorname{tr}(Y_A)Y-Y^2}{\Delta}
                                                       \in\mathcal C.
\]
Moreover
\[
 \overline{H_{00}-H_{22}}
       =(\bar\chi_2(g_*)-1)(\bar\chi_1(g_*)-1)\ne0.
\]
Consequently
\[
 E_{00}=\frac{E_AHE_A-H_{22}E_A}{H_{00}-H_{22}}\in\mathcal C,
 \qquad E_{22}=E_A-E_{00}\in\mathcal C .
 \tag{TA2}
\]
Their adjoints give \(E_{33},E_{11}\).
All four coordinate projectors thus belong to the integral
span, with unit denominators only.

Put \(i\leadsto j\) when \(r(g)_{ij}\ne0\) for some \(g\).
Over the fraction field, (TA2) gives
\[
 i\leadsto j\quad\Longleftrightarrow\quad
       E_{ij}\in\operatorname{span}\{r(g)\}.
 \tag{TA3}
\]
Indeed sandwich by \(E_{ii},E_{jj}\), and divide by the nonzero
entry. Hence this relation is reflexive and transitive,
and symplectic adjoint gives
\(i\leadsto j\Rightarrow 3-j\leadsto3-i\).
The representation is absolutely irreducible precisely when
all sixteen relations hold. This is the matrix-algebra form
of irreducibility, which also avoids any semisimplicity
assumption in positive characteristic.

\paragraph{An isolated odd two-dimensional block is impossible.}
Suppose \(i,j\) lie in different sectors, mutually satisfy
\(\leadsto\), and form a strongly connected component by
themselves. Every multiplication term through another vertex
vanishes: such a term would put that vertex in the same
component. In particular
\[
 r(gh)_{ij}=r(g)_{ii}r(h)_{ij}
                         +r(g)_{ij}r(h)_{jj}.            \tag{TA4}
\]
Use a complete DVR model and choose a nonzero entry
\(b=r(g_1)_{ij}\) of smallest valuation.
Write \(r(g)_{ij}=b\,c(g)\).
Then \(c(g)\) is integral, its residue is locally constant,
and \(c(g_1)=1\).
Equation (TA4) makes
\[
 z(g)=\frac{\overline{c(g)}}{\bar\chi_j(g)}
\]
a cocycle for the odd root
\(\gamma=\bar\chi_i/\bar\chi_j\).

The ordinary local flag induces a flag on the corresponding
two-dimensional subquotient. One way to see this without
assuming a direct-sum decomposition is to order the strongly
connected components, take the invariant coordinate-subspace
filtration, and intersect and project the ordinary flag.
The resulting nonzero rank-two subquotient has an invariant
line. Write its coordinates as \((v,1)\) after interchanging
\(i,j\) if necessary; the valuation-ring property permits
\(v\in\mathcal V\).
If this interchanges the two indices, reselect \(b\) as a
smallest-valuation nonzero entry in that orientation; mutual
reachability supplies it, and the preceding cocycle calculation
applies to the reversed root as well.
Its eigenline equation is
\[
 b\,c(s)=\bigl(r(s)_{ji}v+r(s)_{jj}-r(s)_{ii}\bigr)v
                  \qquad(s\in G_{\mathbb Q_\ell}).        \tag{TA5}
\]
At the actual separator the bracket is a unit: the opposite
off-diagonal entry is in the maximal ideal, and the residual
diagonal values in different sectors are distinct.
Thus \(v/b\) is integral.
Divide (TA5) by \(b\) and reduce. It says that \(z\) is a
coboundary on the full local group at \(\ell\).
In the principal Case II packet \(S_b=\varnothing\), because
\(\psi\) has \(2\)-power order. At every \(p\ne\ell\) the
actual relation for \(r\) therefore kills
\(I_p\cap\ker\bar r\).
This subgroup has index prime to \(\ell\) in
\(I_p\cap\ker\gamma\): the residual representation is a
sum of finite prime-to-\(\ell\) characters.
Thus \(z\) satisfies exactly the unramified-up-to-prime-to-\(\ell\)
index condition in the residual vanishing theorem.
No further bad-prime local condition is needed.

Apply the already proved residual odd-root vanishing
(Proposition~\ref{prop:sr-actual-corx}, equivalently
the six root instances of the Bernoulli-unit calculation).
It follows that
\[
 \overline{c(g)}
           =m\bigl(\bar\chi_i(g)-\bar\chi_j(g)\bigr)
                         \quad\text{for every }g.         \tag{TA6}
\]
The same conclusion over an extended residue field follows
coordinate by coordinate: apply each linear functional over
the original finite residue field to \(z\).
For a nontrivial \(\gamma\), one element with
\(\gamma(g)\ne1\) determines the coboundary scalar; the
coordinate conclusions therefore combine into (TA6).

But (TA2) writes \(E_{ii}=\sum_g a_g r(g)\) with finitely many
integral coefficients. Its \(ij\)-entry, divided by \(b\),
has zero residue. Equation (TA6) makes this residue
\[
 m\sum_g\bar a_g
            \bigl(\bar\chi_i(g)-\bar\chi_j(g)\bigr)=m,
\]
using the \(ii\) and \(jj\) entries of that same projector
identity. Hence \(m=0\).
This contradicts \(c(g_1)=1\).
There is therefore no isolated strongly connected pair
joining different sectors.

\begin{srlemma}\label{lem:sr-four-actual-triangular-alternatives}
An actual normalized rank-four point has one of the following
three forms:
\[
 \begin{array}{ll}
 \mathrm{(I)}&r\text{ is absolutely irreducible};\\
 \mathrm{(D)}&
  i\ne j,\ i\leadsto j\ \Longrightarrow\ j\not\leadsto i;\\
 \mathrm{(B)}&
 0\leadsto2,\ 2\leadsto0,\text{ and all cross-sector entries}\\
 &\text{in one of the two block directions vanish}.
 \end{array}
\]
In the last case both diagonal blocks are absolutely
irreducible.
\end{srlemma}
\begin{proof}
Strongly connected components are carried to strongly connected
components by \(i\mapsto3-i\).
A component of size three would meet its image and thus equal
it, impossible for a fixed-point-free involution on three
vertices. A component of size four gives (I).
A component of size two joining different sectors was just
excluded. The only remaining pair is \(A=\{0,2\}\), whose
image is \(B=\{1,3\}\).
These are both components, and their quotient directed graph
has no cycle; hence entries in at least one direction vanish.
Equation (TA3) gives the full matrix algebra on each
two-dimensional diagonal block.
If there is no nonsingleton component, (D) holds.
\end{proof}

\paragraph{Take the block-diagonal part without losing ordinarity.}
In case (B), erase all cross-sector entries and call the result
\(r_0\). A block-triangular product shows that \(r_0\) is a
representation, with the same characteristic polynomials.
Symplectic adjoint commutes with this operation and preserves
its multiplier.

The ordinary flag is still valid.
The matrix \(r(\varphi)\) in (TA1) is already block diagonal
and preserves that flag. The polynomial formula producing
\(E_A\) shows that \(E_A\) also preserves it.
In a flag-adapted basis both \(r(s)\) and \(E_A\) are upper
triangular; therefore
\[
             r_0(s)=E_A r(s)E_A+(1-E_A)r(s)(1-E_A)
\]
is upper triangular with exactly the same diagonal.
This preserves both ordered residual local characters and
the exact prescribed inertial characters.
Erasing entries preserves every killed-inertia relation and
the complete DVR model.
Thus \(r_0\) is an actual block-diagonal ordinary point to
which Proposition~\ref{prop:sr-four-actual-block-extension}
applies.

\paragraph{The diagonal alternative in characteristic \(\ell\).}
Assume (D), characteristic \(\ell\), and the original
fixed-inertia weight-zero specialization:
the ordinary character on each \(I_\ell\) is exactly the
image of the prescribed \(\chi_i|_{I_\ell}\), whose
characteristic-\(\ell\) image is its finite residual character.
No unspecialized ordinary weight variable is allowed in this
paragraph. The curve used in the finite-weight-zero argument
below has this condition by construction.
The diagonal entries \(\delta_i(g)=r(g)_{ii}\) are characters:
every off-diagonal term in a diagonal entry of \(r(g)r(h)\)
would give a two-cycle.

We check that these characters have finite order.
Let \(U=\ker\bar r\cap\ker\psi\), an open normal subgroup.
At inertia at \(\ell\), an element of \(U\) has all ordinary
diagonal characters equal to \(1\); the representation is
unipotent there. A strictly upper \(4\) by \(4\) matrix has
fourth power zero, so in characteristic \(\ell\)
\[
                  r(s)^{\ell^2}=1
                   \quad(s\in U\cap I_\ell),             \tag{TA7}
\]
including \(\ell=3\).
At every other prescribed prime the actual killed-inertia
condition applies on \(U\); elsewhere the global unramified
condition applies.
Thus, with \(m=[G_{\mathbb Q}:U]\), every
\(\delta_i^{m\ell^2}\) is unramified at every finite prime.
At every finite DVR precision its image is finite.
Minkowski's theorem says that the corresponding everywhere
unramified finite extension of \(\mathbb Q\) is trivial.
Taking all precisions yields
\[
                         \delta_i^{m\ell^2}=1.            \tag{TA8}
\]
These characters therefore take values algebraic over
\(\mathbb F_\ell\).

There is a linear functional on the root lattice strictly
positive on every occurring nonzero off-diagonal weight.
For an explicit formula put
\[
 a_i=\#\{j:i\leadsto j\},\qquad h_i=a_i-a_{3-i},\qquad
 L(x,y)=(h_0-h_1)x+(h_0-h_2)y .
\]
If \(i\ne j\) and \(i\leadsto j\), transitivity and absence of
a two-cycle give \(a_i>a_j\); applying the reflected relation
gives \(a_{3-j}>a_{3-i}\).
Thus
\[
                  L(w_{ij})=h_i-h_j>0.                  \tag{TA9}
\]
Consequently a nonzero balanced entry monomial contains no
off-diagonal factor. It is a product of diagonal character
values, and (TA8) makes it algebraic over \(\mathbb F_\ell\).
Every balanced monomial \(x\) therefore satisfies
\(x^{\ell^n}=x\) for some \(n>0\).
This proves the required constant-character conclusion,
rather than just a triangularity assertion.

The source calculations expanded here are
\srcid{SymplecticGSp4.EiiUnit}, \srcid{GeAlgTop},
\srcid{CompactIntegralFields.CobOfCoord.triV},
\srcid{BlockPart}, and \srcid{DiagConst}.
